% TOP_PROBLEM: {"id": 2856, "title": "Kirby Problem 3.58", "category_id": 7, "category": {"id": 7, "name": "topology", "display_name": "Topology", "description": "Properties preserved under continuous deformations.", "slug": "topology", "order_index": 7, "created_at": "2026-07-31T15:26:25.671Z"}, "status": "open"}
% FIRST_POSTED: null
% ENTRY_KIND: statement_only
% Imported from: batch12.tex; UnsolvedMath ID 2856
\documentclass[11pt]{article}
\usepackage[margin=25mm]{geometry}
\usepackage{fontspec}
\setmainfont{Latin Modern Roman}
\usepackage{amsmath,amssymb,mathtools}
\usepackage{unicode-math}
\setmathfont{Latin Modern Math}
\usepackage{xurl,hyperref}
\hypersetup{colorlinks=true,urlcolor=blue}
\setcounter{secnumdepth}{0}
\setlength{\parindent}{0pt}
\setlength{\parskip}{5pt}
\setlength{\emergencystretch}{3em}
\newcommand{\sref}[2]{\hyperlink{source:#1:#2}{[S#2]}}
\newcommand{\cataloguescope}[1]{\paragraph{Recorded target.}#1}
\begin{document}
% UnsolvedMath ID 2856; source catalogue code KP-3.58
\setcounter{section}{2855}
\section{Kirby Problem 3.58}\label{2856}
{\small\sffamily UnsolvedMath ID 2856 · Topology}
\subsection{Definitions and mathematical statement}

\paragraph{Shared notation.}$\mathbb N=\{1,2,\ldots\}$, and $\mathbb Z,\mathbb Q,\mathbb R,\mathbb C$ denote the integers, rationals, reals and complex numbers. Any other conventions are specified locally.


For a closed oriented $3$-manifold $Y$ and a spin$^c$ structure $\mathfrak s$, write $\operatorname{HF}^{\circ}(Y,\mathfrak s)$ for Heegaard Floer homology and $\operatorname{HM}^{\circ}(Y,\mathfrak s)$ for its corresponding monopole Floer version. For $\circ=-,+,\infty$, the monopole versions are respectively the hat, check and bar monopole theories; the hat Heegaard theory corresponds to the associated $U$-mapping-cone version. Use matching coefficients, gradings, $U$ actions and orientation conventions in both theories, with canonical coherent orientations for integral Heegaard coefficients.
A spin$^c$ cobordism $(W,\mathfrak t):(Y_0,\mathfrak s_0)\to(Y_1,\mathfrak s_1)$ is a compact oriented $4$-manifold with boundary $-Y_0\sqcup Y_1$ and a spin$^c$ structure restricting to $\mathfrak s_i$. Its Floer maps are $F^{\circ}_{W,\mathfrak t}$ and $\operatorname{HM}^{\circ}_{W,\mathfrak t}$, using matching homology orientations. The mixed Ozsv\'ath--Szab\'o invariant $\Phi_{X,\mathfrak t}$ of a closed $4$-manifold is the mixed minus-to-plus Floer cobordism invariant, while $\operatorname{SW}_{X,\mathfrak t}$ counts the corresponding Seiberg--Witten solutions. Their comparison uses the same homology orientation, insertions and, whenever necessary, chamber data, in every setting where these invariants are defined.
\paragraph{Statement recorded in the supplied source.}
Construct isomorphisms $\Psi^{\circ}_{Y,\mathfrak s}$ between the corresponding Heegaard and monopole Floer groups, compatible with their structures and satisfying, for every spin$^c$ cobordism,
\[
\Psi^{\circ}_{Y_1,\mathfrak s_1}\circ F^{\circ}_{W,\mathfrak t}
=\operatorname{HM}^{\circ}_{W,\mathfrak t}\circ\Psi^{\circ}_{Y_0,\mathfrak s_0},
\]
with the standard common sign convention over $\mathbb Z$. In particular, prove equality of $\Phi_{X,\mathfrak t}$ and $\operatorname{SW}_{X,\mathfrak t}$ for closed $4$-manifolds, using corresponding orientation and chamber conventions on their common domains of definition. \sref{2856}{2}
\subsection{Short English statement}
Can the isomorphisms between Heegaard and monopole Floer homology be chosen to respect every cobordism map, so that the closed four-manifold mixed invariant equals the Seiberg--Witten invariant?
\subsection{Sources}
\begin{description}
\item[\hypertarget{source:2856:1}{S1}] ULAM.AI and the UnsolvedMath Contributors, \emph{UnsolvedMath: A Curated Collection of Open Mathematics Problems}, record 2856 (KP-3.58). CC BY 4.0. \url{https://huggingface.co/datasets/ulamai/UnsolvedMath}.
\item[\hypertarget{source:2856:2}{S2}] R. İnanç Baykur, Robion C. Kirby and Daniel Ruberman (editors), \emph{K3: A New Problem List in Low-Dimensional Topology}, Mathematical Surveys and Monographs 295 (2026), Problem 3.58, p. 173 and following remarks; original author version. \url{https://math.berkeley.edu/sites/default/files/surv-295-ruberman-watermarked-author-pdf.pdf}.
\end{description}
\label{end:2856}
\end{document}
